\section{BIBLIOGRAFÍA Y REFERENCIAS A CONSULTAR}

Para la fundamentación técnica y jurídica del presente dictamen, el \textbf{Equipo Pericial \GrupoPracticas} se ha basado, además de en su formación y experiencia profesional, en la normativa legal, estándares técnicos, bibliografía especializada y documentación oficial que se detallan a continuación:

\textbf{Marco normativo y legislación procesal}
\begin{itemize}[leftmargin=1.5em, itemsep=0.35em]
	\item \textbf{Ley 1/2000, de 7 de enero, de Enjuiciamiento Civil (LEC).} Boletín Oficial del Estado, núm. 7, de 8 de enero de 2000. Artículos 335.2 (juramento o promesa de objetividad del perito), 340.1 (condiciones de titulación oficial habilitante) y 343 (tachas de los peritos).
\end{itemize}

\textbf{Normas técnicas de informática forense y seguridad de la información}
\begin{itemize}[leftmargin=1.5em, itemsep=0.35em]
	\item \textbf{UNE-EN ISO/IEC 27037:2016.} \textit{Tecnología de la información. Técnicas de seguridad. Directrices para la identificación, recogida, adquisición y preservación de evidencias digitales.} Asociación Española de Normalización (UNE).
	\item \textbf{UNE 71506:2013.} \textit{Tecnologías de la Información (TI). Metodología para el análisis forense de las evidencias electrónicas.} Asociación Española de Normalización (UNE).
	\item \textbf{UNE-EN ISO/IEC 27001.} \textit{Seguridad de la información, ciberseguridad y protección de la privacidad. Sistemas de gestión de la seguridad de la información.} Asociación Española de Normalización (UNE).
\end{itemize}

\textbf{Estándares de comunicaciones electrónicas (IETF RFC)}
\begin{itemize}[leftmargin=1.5em, itemsep=0.35em]
	\item \textbf{Klensin, J. (2008).} \textit{RFC 5321: Simple Mail Transfer Protocol (SMTP).} Internet Engineering Task Force (IETF). Disponible en: \url{https://www.rfc-editor.org/info/rfc5321}.
	\item \textbf{Resnick, P. (2008).} \textit{RFC 5322: Internet Message Format.} Internet Engineering Task Force (IETF). Disponible en: \url{https://www.rfc-editor.org/info/rfc5322}.
	\item \textbf{Kitterman, S. (2014).} \textit{RFC 7208: Sender Policy Framework (SPF) for Authorizing Use of Domains in Email.} Internet Engineering Task Force (IETF).
	\item \textbf{Crocker, D., Hansen, T., \& Kucherawy, M. (2011).} \textit{RFC 6376: DomainKeys Identified Mail (DKIM) Signatures.} Internet Engineering Task Force (IETF).
	\item \textbf{Kucherawy, M., \& Zwicky, E. (2015).} \textit{RFC 7489: Domain-based Message Authentication, Reporting, and Conformance (DMARC).} Internet Engineering Task Force (IETF).
\end{itemize}

\textbf{Bibliografía técnica y documentación oficial de herramientas}
\begin{itemize}[leftmargin=1.5em, itemsep=0.35em]
	\item \textbf{Carrier, B. (2005).} \textit{File System Forensic Analysis.} Addison-Wesley Professional. (Análisis estructural del sistema de archivos NTFS, Tabla Maestra de Archivos \texttt{\$MFT}, atributos residentes y flujos de datos alternativos).
	\item \textbf{Russinovich, M., Solomon, D. A., \& Ionescu, A. (2017).} \textit{Windows Internals, Part 1 \& Part 2.} Microsoft Press. (Arquitectura interna de sistemas Windows, registro de eventos \texttt{.evtx}, identificadores de seguridad \texttt{SID} y claves del Registro).
	\item \textbf{Carrier, B.} \textit{The Sleuth Kit (TSK) \& Autopsy Digital Forensics Documentation (v. 4.23.1).} Disponible en: \url{https://www.sleuthkit.org/}.
	\item \textbf{Exterro Inc.} \textit{FTK Imager User Guide (v. 8.3.0.27).} Documentación técnica de adquisición forense en modo solo lectura y verificación criptográfica.
	\item \textbf{Microsoft Corporation.} \textit{Documentación técnica oficial de Windows Server 2008 R2 y Windows 11} (utilidades \texttt{certutil}, \texttt{Get-FileHash}, administración de recursos compartidos SMB mediante \texttt{net share} y directivas de almacenamiento \texttt{StorageDevicePolicies}). Disponible en: \url{https://learn.microsoft.com/}.
\end{itemize}